\chapter{Optimisation of Thunks in CakeML}
\label{chap:thunks}

Moving downwards on the compilation pipeline, this chapter describes our work
for the last part of the thesis, optimising the representation of thunks in
CakeML. The CakeML compiler was not built for a lazy language, so it rightly
doesn't consider thunks natively. However, when considered as a target for
PureCake, it is critical that it can handle thunks effectively, since they are
the main source of inefficiency between lazy and strict programs. In our work,
we add thunks to the CakeML source language's state, and target them directly
from PureCake. We then consider the consequences of this change on both
codebases and discuss possible optimisations in CakeML's runtime.

\section{Strengthening the Semantics of \texttt{force}}

As discussed in \autoref{chap:intro}, \ThunkLang{} is used as a phase to change
the evaluation strategy from call-by-name to call-by-value with the addition of
thunks and operations for thunks to simulate laziness. The semantics for the
thunk operations are quite simple (we ignore clocking for the sake of clarity):

\begin{figure}[h]
  \[
    \inference{}{|[\lit{delay}\ e|] = \lit{Thunk}\ e} \qquad
    \inference
    {|[x|] = \lit{Thunk}\ e& |[e|] = u}
    {|[\lit{force}\ x|] = u}
  \]
  \caption{Simplified derivation rules of \ThunkLang{}'s special operators}
\end{figure}

A \lit{delay} boxes an expression into a thunk and \lit{force} queries
a thunk for its expression and evaluates it. However, we would like to have a
stronger version of \lit{force} which asserts that the result is not itself a
thunk. This provides a very useful invariant that we use in proofs throughout
the compilation pipeline. Therefore, we change the semantics of \lit{force} to
include that check.

\begin{figure}[h]
  \[
    \inference
    {|[x|] = \lit{Thunk}\ e& |[e|] = u& \nexists t. u = \lit{Thunk}\ t}
    {|[\lit{force}\ x|] = u}
  \]
  \caption{Stronger version of \lit{force}}
\end{figure}

To see how this small change affects compilation, we need to inspect the
\ThunkLang{} phase of the compiler closer.

\section{Adapting \ThunkLang{} to the New Semantics}

Before the PureCake compiler transforms a \ThunkLang{} program into the next
intermediate representation \EnvLang{}, it first goes through a number of
\ThunkLang{}-to-\ThunkLang{} transformations that perform optimisations on the
program. For instance, the simple transformation \lit{force\_rel} removes
\lit{force}s around variables that are known to already be \lit{force}d.
This is useful when we have an expression like \hsinl{let m = force n in ...},
where any occurrence of \lit{force n} inside the body is rewritten to just
refer to \lit{m}. The transformation that concerns us the most is the
\lit{delay\_force} pass, which transforms expressions of the form
\lit{delay (force ($e$))} to just $e$. This transformation in its current form
cannot be proved correct with our new assertion. To explain why, we must examine
how exactly we prove the correctness of these optimisation passes.

\subsection{Proving correctness of \ThunkLang{} Optimisations}

Formally, a transformation is a pair of binary relations $(R_e,R_v)$. The
relation $R_e$ describes the transformation of expressions, and $R_v$ the
appropriate transformation of the values resulting from evaluating a transformed
expression. Therefore, the main theorem to prove for each transformation is that
evaluating related expressions results to related values.

\begin{theorem}[\lit{exp\_rel\_eval}] For a transformation $(R_e,R_v)$ and
  any expressions $e_1,e_2$, if $e_1$ and $e_2$ are related under $R_e$ and
  $eval(e_1)$ is not an error,
  then the results of $eval(e_1)$ and $eval(e_2)$ must be related under $R_v$.
\end{theorem}

Then, the whole process of translating a \ThunkLang{} program to an \EnvLang{}
program consists of a chain of transformations $(R_{e_1},R_{v_1}), \cdots,
(R_{e_n},R_{v_n})$ such that each one satisfies \lit{exp\_rel\_eval} and that
for any \ThunkLang{} expression $e$, there exist \ThunkLang{} expressions
$e_1,\cdots,e_n$ such that $R_{e_1}(e,e_1) \land R_{e_2}(e_1,e_2) \land \cdots
\land R_{e_n}(e_{n-1},e_n)$. Finally, a transformation $R_{te}(e_n,e_e)$ exists
that transforms the last \ThunkLang{} expression to an \EnvLang{} expression.

On their own, these relations are not strong enough to prove program
equivalence, they only prove a relation between the starting and the final
program. To prove equivalence from these relations, we show that they are
bisimulations \cite{howe1996proving}. A binary relation is a bisimulation if its
two parts behave the same according to some rules. Our relations are over
\ThunkLang{} expressions and the rules are the evaluation function, and using
this theory we can prove that the evaluation of two related expressions cannot
be distinguished, therefore making them equivalent.

It is evident that each expression relation must consider every expression on
the left hand side that can be generated by its previous relation. Therefore,
each relation has boilerplate cases that just carry over each expression
unchanged, and some interesting once that apply the transformations. For
example, the \lit{delay\_force} relation has two interesting cases
\[
\inference{v\ \text{variable}}{R_e(\lit{delay (force ($v$))},v)}
\quad
\inference{R_v(v_1,v_2)}{R_e(\lit{delay (force ($v_1$))},v_2)}
\]
but also all the trivial
cases such as \[\inference{
  R_e(x_1, x_2)&
  R_e(y_1, y_2)&
  R_e(z_1, z_2)}
{R_e(\lit{if } x_1 \lit{ then } y_1 \lit{ else } z_1,
\lit{if } x_2 \lit{ then } y_2 \lit{ else } z_2)}.\]

The reason we allow the first expression in each relation to be an error is
because we only care not to generate a failing program from a correct one, and
not the other way around. In the case of \lit{delay\_force}, for a variable $v$
that expands to a thunk that fails when forced, the program \lit{delay (force
($v$))} will fail, however the transformed program which is just $v$ may not
fail if the result of expanding $v$ is never forced.

Since we added this new assertion in the semantics of \lit{force}, we need to
prove that none of the optimisations violates it. For each of these
transformations, our change in the semantics now requires a new theorem to be
proved:

\begin{theorem}[\lit{v\_rel\_thunk}] For a transformation $(R_e,R_v)$ and any
values $v_1,v_2$, if $v_1$ and $v_2$ are related under $R_v$, then $v_1$ is a
thunk value if and only if $v_2$ is also a thunk value
\end{theorem}

For all the transformations implemented this was trivial to prove, except for
\lit{delay\_force}. To see why, we consider what happens when we try to prove
the evaluation theorem for the interesting case of $Re(\lit{delay (force
($v_1$))}, v_2)$ when $R_v(v_1,v_2)$: The expansion of $eval(\lit{delay (force
($v_1$))})$ will eventually assert that the result of the inner \lit{force} is
not a thunk. However, for $v_2$ no such assertion exists, therefore
\textbf{\lit{v\_rel\_thunk}} is not true and the evaluation proof fails. To
explain how we correct \lit{delay\_force}, we first need to briefly dive into
how PureCake handles side effects.

\subsection {Side Effects in a Pure Language}

So far we have only concerned ourselves with the pure semantics of \ThunkLang{}.
That is, the semantics of expressions that are perfectly
referentially transparent and don't rely on any implicit state. This is the core
of pure functional programming, but in any programming language there needs to
be a way to interact with the outside world. Any operation that relies on such a
state is called a side effect, and examples of side effects include user input
and output, global state, exception handling and many others.

The way pure languages handle side effects is by introducing monads to the
language \cite{benton2000monads}. A monad in the context of programming
languages is a wrapper around a computation which also produces some additional
information about that computation, like an execution trace or non-determinism.
PureCake also employs monads to handle side effects. The evaluation of a
\ThunkLang{} program at the top level looks like a loop over monadic actions
that might possibly interact with state, such as writing and reading from a
mutable array or writing to the console. Internally, these operations call the
usual pure semantics that we have discussed in depth without exposing any of the
state that is being handled.

To be able to reason about stateful computation, PureCake uses Interaction Tree
Semantics \cite{xia2019interaction}. Interaction trees are a data structure for
representing the behaviour of programs that interact with their environments.
Interaction trees support the definition of interpreters that handle actions
using monads, and bisimulation proofs can be used to show equivalence between
interaction trees, similarly with how we described equivalence proofs for
program transformations.

For example, we can consider a simple function that interleaves pure computation
with side effects
\begin{minted}{haskell}
foo :: Int -> IO Int
foo x = do
  let y = f x
  print y
  z <- input
  return z
\end{minted}
When this function is ran, the \lit{let} expression will be evaluated first. Its
result will then be printed to the terminal using the stateful \lit{print}
function. Then, \lit{input} is used to read an integer from the terminal and
unwraps it to the pure value \lit{z}, which is then wrapped in a monadic value
by \lit{return} so it matches the return type of the function. This
back-and-forth between pure and impure evaluation is modelled by the IO monad,
which stands for Input/Output.

Because the results of stateful computations can be used by regular functions,
they must be wrapped in thunks. Therefore, any value acquired by \lit{input} or
\lit{return} or any other returning monadic action must be delayed. However,
these monadic actions can also be called and chained by other monadic actions,
where it's not necessary to wrap the results in thunks. Thus, in \ThunkLang{}
the semantics of the monadic operators themselves are strict, and there is a
step during the translation from PureCake's source language to \ThunkLang{} that
wraps the program's monadic actions that need thunks with \lit{delay}s. For
example, whenever there is a call to \lit{input} in the source language, in
\ThunkLang{} it is translated to \hsinl{input >>= \v -> return (delay v)},
meaning that instead of directly returning the result of \lit{input}, we first
wrap it in a thunk since the monadic interpreter does not do this on its own.

This optimisation is crucial for avoiding the allocation of unnecessary thunks
during monadic computations, but complicates the semantics of the language and
is the root of our problem with \lit{delay\_force} as we will now see.

\subsection {Fixing \lit{delay\_force}}

We saw that the interesting cases in \lit{delay\_force} are concerned with value
and variable expressions. As mentioned in \autoref{chap:intro}, \ThunkLang{}
uses a substitution based semantics. During \ThunkLang{}'s evaluation,
encountering a variable expression immediately results in an error. Instead,
before evaluating an expression, every free variable in its body is substituted
for some value expression. A value expression is just the embedding of a value
inside the program, so when the evaluation encounters it, it just returns the
value as is.

We can fix the failing proof by making two assumptions:
\begin{itemize}
\item Every value expression surrounded by a \lit{force} is a thunk
\item Every variable expression is substituted with a value expression of a
  thunk
\end{itemize}
These might seem like strong assumptions, but we just have to look at how the
compiler transitions from PureCake's source language to \ThunkLang{} to see that
this holds, except for one place, the handling of the monadic side effects.
Since the compiler optimises away certain \lit{delay}s in the monadic
interpreter, this invariant doesn't hold.

\begin{figure}[h]
  \begin{tikzpicture}
    \tikzstyle{box} = [
      rectangle, rounded corners, minimum width=2cm, minimum height=1cm,
      text centered, draw=black
    ]
    \node (box1) [box] {PureCake source};
    \node (box2) [box, right=0.7cm of box1] {
      \shortstack {
        \texttt{ThunkLang} \\ Optimised monadic interpreter
      }
    };
    \node (box3) [rectangle, text centered, right=0.7cm of box2] {
        \shortstack {\lit{delay\_force} and rest \\ of optimisation pipeline}
    };
  \draw[->](box1) to [bend left,looseness=1] node[above]
    {\shortstack{Wrap monadic results in \\
     the program with \lit{delay}s}} (box2);
  \draw[-stealth] (box2.east) -- (box3.west) node[]{};
  \end{tikzpicture}
  \caption{Current translation of monadic expressions from PureCake source to
           \ThunkLang{}}
\end{figure}

To mitigate this issue we carried out a 4-step plan:
\begin{itemize}
  \item Implement an alternative unoptimised monadic interpreter that adds
    \lit{delay}s around the results of all the monadic actions
  \item Remove the wrapping of monadic results during the PureCake source to
    \ThunkLang{} pass to avoid duplicated \lit{delay}s
  \item Implement a transformation from the unoptimised monadic interpreter to
    the optimised monadic interpreter by adding back the wrapping of monadic
    results in \lit{delay}s
  \item Move \lit{delay\_force} before the transformation to the optimised
    interpreter happens
\end{itemize}

\begin{figure}[h]
  \begin{tikzpicture}
    \tikzstyle{box} = [
      rectangle, rounded corners, minimum width=2cm, minimum height=1cm,
      text centered, draw=black
    ]
    \node (box1) [box] {PureCake source};
    \node (box2) [box, right=1.5cm of box1] {
      \shortstack {
        \texttt{ThunkLang} \\ Unoptimised \\ monadic interpreter
      }
    };
    \node (box3) [box, right=1.5cm of box2] {
      \shortstack {
        \texttt{ThunkLang} \\ Unoptimised \\ monadic interpreter
      }
    };
    \node (box4) [box, text centered, below=1.5cm of box3] {
        \shortstack {\ThunkLang{} \\ Optimised \\ monadic interpreter}
    };
    \node (box5) [rectangle, text centered, left=1.5cm of box4] {
        \shortstack {Rest of compilation \\ pipeline}
    };
    \draw[-stealth] (box1.east) -- (box2.west) node[]{};
    \draw[->](box2) to [bend left,looseness=1] node[above]
      {\shortstack{\lit{delay\_force}}} (box3);
    \draw[->](box3.west) to [bend right,looseness=1] node[right]
      {\shortstack{Wrap monadic results in \\ the program with \lit{delay}s}}
      (box4.west);
    \draw[-stealth] (box4.west) -- (box5.east) node[]{};
  \end{tikzpicture}
  \caption{Adapted translation of monadic expressions from PureCake source to
           \ThunkLang{}}
\end{figure}

With this translation scheme, we were able to carry on our invariant that the
result of \lit{force} never returns a thunk without sacrificing any
optimisations. This also serves as a prime example of the usefulness of formal
verification in compiler construction. All this rearranging of fragile parts
would be very difficult to achieve correctly and optimally in a conventional
compiler. Since we implement everything in a theorem prover we can be sure that
we achieved our goal and didn't lose any optimisations in the process.

\section{The Stateful Semantics of \StateLang{}}

The next intermediate representation from \ThunkLang{} is \EnvLang{}. This
language doesn't transform thunks or their operations in any way, so we just had
to add the semantic check after the \lit{force} and everything else stayed the
same. Therefore we can immediately move on to the more interesting \StateLang{}.

\StateLang{} is the last intermediate representation in the PureCake pipeline,
and the boundary between PureCake and CakeML. It is also the only intermediate
representation with a stateful semantics which means it can represent stateful
thunks. Before \StateLang{}, thunks are stateless; when evaluating an expression
like \hsinl{f (g x)}, the argument \lit{g x} is turned into a thunk value and
whenever the corresponding parameter is used in the body of f, the expression
\lit{g x} is computed anew without the possibility of memoising the result, as
that would require some stateful operation.

\StateLang{}'s semantics provide a solution for this problem by employing a
stateful store. This store is an array of mutable arrays, where each array
element can be mutated in place. \StateLang{} also provides primitive operations
for manipulating arrays, namely \lit{Alloc} for allocating an array, \lit{Sub}
for dereferencing an array element and \lit{Update} for updating the value of an
array element.

\StateLang{} uses a continuation-based small-step semantics
\cite{nakata2009small}, in contrast to \ThunkLang{}'s big step semantics. In
small-step semantics, the semantic function does not walk the expression
structure recursively. Instead, it breaks down the evaluation process into
atomic steps, and takes one step at a time, returning the rest of the
computation as a result. Concretely, the evaluation function takes a state, a
clock, and an expression and returns one of three results: An expression, a
value, or an error. In actuality, there are two more results, namely exception
and action values for performing cross-language calls but they are not important
for our treatment of the semantics. Along with the result, the \lit{step}
function returns a new state and a continuation stack. The continuation stack is
a list of remaining operations waiting to be executed in the rest of the steps.
As an example, let's consider how \StateLang{}'s semantics handle application
expressions. The (simplified) \lit{step} function looks like this:

\begin{minted}{haskell}
step st k (Exp env (App op (e:es))) =
  if not (num_args_ok op (length (e:es)))
    then error st k
    else push env e st (AppK env op [] es) k
\end{minted}

We see that first the function checks if the number of arguments is the correct
for the given operation. In \StateLang{}, the term ``application operator'' is
used to refer to many things, one of them being a function application, but for
example another one being the ``application'' of the \lit{Sub} operator to an
index and an array. So, assuming that we are talking about function application
for this example, the \lit{step} function should have received two expressions
else it will return an error. If the number of arguments is correct, the
function calls \lit{push}, which just constructs a result consisting of \lit{e}
as the expression, \lit{st} as the state and the continuation stack \lit{k}
prepended with the \lit{AppK} continuation.

Intuitively, the \lit{step} function takes away arguments from the application
expression one by one, evaluates them and has stored a ``promise'' in the form
of a continuation that it will eventually apply the operator to the resulting
values. When all the operands have been evaluated, the \lit{AppK} continuation
will come to be handled, which will apply the operator to the list of values
produced by the evaluation of the arguments. In the case of the application
operator, it will check the head in the list of values to make sure it is a
closure; then it will bind the second value in the list to the closure's
environment, and finally it will return the closure's expression with the new
environment for the \lit{step} function to start evaluating it.

We can now explain the semantics of the stateful operators \lit{Alloc},
\lit{Sub} and \lit{Update}. These operators are all application operators,
therefore their evaluation will take a similar course with the application
operator. When applied to the list of values, \lit{Alloc} expects two values: a
number $n$ and an arbitrary value $v$. The \lit{application} function then will
append a new array of length $n$ to the state and fill it with $n$ occurrences
of $v$. It will then return the length of the state as the index of the
allocated array, assuming that accesses to the state are zero-based. The
\lit{Sub} operator also expects two values: an index $n$ to select an array from
the state, and an index $m$ to select an element from the $n$th array of the
state. The \lit{application} function returns the selected element and leave the
state intact. Lastly, the \lit{Update} operator expects three values: two
indices to select an element like \lit{Sub}, and a third value to replace that
element. The \lit{application} function will update that element in the state
with the third value, and return the unit constructor.

Currently, \StateLang{}'s store is used to implement thunks in an ad-hoc way.
Instead of having a proper representation for thunks, a thunk is represented as
a two value array: the first value is a boolean flag and the second value is the
actual value. When a thunk is needed, the array is allocated with the flag set
to signal that the thunk is not yet evaluated, and the second value being a
closure that when applied to unit computes the thunk's value. When a \lit{force}
encounters such a value, it checks the flag in the first place of the array and,
if it is set it immediately returns the value in the second place. If it is not
set, it applies the second value to unit, replaces the function with the value,
sets the flag and finally returns the value. To make these operations more
efficient, there are dual unsafe operators of \lit{Sub} and \lit{Update} called
\lit{UnsafeSub} and \lit{UnsafeUpdate} respectively. Since the semantics know
the shape the array argument of a \lit{force} will have, there is no need to
check if the two indices are in range.

The transition from the stateless thunks of \ThunkLang{} to the stateful thunks
of \StateLang{} is done in a \StateLang{}-to-\StateLang{} transformation; the
\StateLang{} source language only has stateless thunks with the usual
\lit{force} and \lit{delay} operators, but implements an optimisation pass that
replaces the stateless thunk expressions with stateful operations on the store,
an optimisation called unthunking.

The decision to design \StateLang{}'s store like that was not arbitrary. This is
the form of CakeML's state and, since \StateLang{} is compiled directly to
CakeML it was the natural choice to make. However, it is evident that this
representation is not ideal and throws away information that can be used during
the CakeML compilation to optimise the program even more.

\section{Making Thunks Explicit}

Before being able to make any optimisations to thunks on the CakeML side, we
first need to make thunks explicit in both the CakeML's source language and in
\StateLang{}.

On the side of CakeML, adding thunks to the language consisted of creating a new
kind of store besides the generic arrays to specialise for thunks, along with
operations for allocating, updating and forcing them. Now, the PureCake compiler
can target them directly instead of simulating them with arrays of two values.
The same change was made in \StateLang{}, with the store now consisting of an
array of arrays or thunks, and with the unsafe application operators removed in
favor of new operators \lit{AllocMutThunk}, \lit{UpdateMutThunk}, and
\lit{ForceMutThunk}. These changes alone made a big difference in the handling
of thunks; compilation from PureCake to CakeML is more straightforward and
obvious, and proofs inside \StateLang{} become easier to develop and understand.
For example, the unthunking optimisation in \StateLang{} previously had no
direct way of forcing a thunk simulated by an array. Instead, it generated code
that looked like this:

\begin{minted}{haskell}
Let "v1" (App UnsafeSub [Var "v", Int 0]) $
Let "v2" (App UnsafeSub [Var "v", Int 1]) $
  If (Var "v1") (Var "v2") $
    Let "wh" (App AppOp (Var "v2") Unit) $
    Let "_" (App UnsafeUpdate [Var "v"; Int 0; True]) $
    Let "_" (App UnsafeUpdate [Var "v"; Int 1; Var "wh"]) $
      Var "wh"
\end{minted}

This code mimics the desired semantics for thunks. For a variable $v$ that
points to a simulated thunk in the store, extract the flag and the value and if
the flag is set then return the value directly, else compute it, update the
thunk and return the value. This means that the semantics of thunks were hidden
as a code generation in a proof instead of being explicit in the semantics of
the language. With the addition of thunk to the language, this whole construct
has just been replaced with a $\lambda$\hsinl{x. App ForceMutThunk [x]}.

By making the semantics more explicit, we are also forced to handle edge cases
of our semantics explicitly. For example, because of the small-step semantics,
when we force a thunk we check if it's evaluated or not. If it isn't, we can't
just evaluate the underlying closure. Instead, we return a configuration that
describes this computation. When we come to handle this continuation at a later
step, we again find the same thunk in the store and try to evaluate it. But we
now need a proof that, in the meantime, this thunk has not already been
evaluated, or is not in the process of being evaluated. This is reminiscent of
the black hole problem in Haskell \cite{van1999optimising}. In Haskell, the
runtime representation of thunks does not have a binary flag, evaluated and not
evaluated, as we have described for our thunks. Instead, it has three possible
values: evaluated, not evaluated and evaluating. When a thunk starts getting
evaluated, the evaluating flag is set. Then, if a \lit{force} tries to update an
evaluating thunk, then the runtime system detects that as an infinite loop and
throws an exception. However, in our semantics we prove that if we try to force
an evaluating thunk we will eventually diverge, so we only have to handle the
evaluating and not evaluating cases in the rest of our proofs and in the
runtime.

\section{Taking Advantage of Thunks in the Runtime}

CakeML, like every high-level functional language, has a garbage collected
runtime. Each time a certain threshold of allocated memory is surpassed, the
program stops running and the garbage collector traverses the memory starting
from known used locations called roots, copying the reachable data into a new
heap and disposing of the old one. The details of garbage collection in CakeML
are described in \cite{sandberg2019verified}, but this mental model is enough
for our description here.

By having thunks represented explicitly in the CakeML source language, we can
make them explicit through all of CakeML's intermediate representations down to
the runtime. This would give the ability to the garbage collector to
discriminate between chunks of memory that represent thunks, which would allow
an important optimisation called updating \cite{seward1992generational}.

\begin{figure}[h]
  \centering
  \begin{tikzpicture}
    \tikzstyle{box} = [rectangle, text centered, draw=black]
    \node (box0) [box, draw=none] {$x$ : };
    \node (box1) [box, right=0cm of box0] { $ptr$ };
    \node (box2) [box, right=1cm of box1, draw=none] { $ptr$ : };
    \node (box3) [box, right=0cm of box2] { \lit{Evaluated}, 5 };
  \end{tikzpicture}
  \caption{Evaluated Thunk Before Optimisation}
  \label{fig:thunk}
\end{figure}

To illustrate the idea of the optimisation, suppose we have a thunk that has
evaluated to the value 5 stored in variable location $x$. The memory block at
$x$ will hold a pointer to the thunk block, which will hold the actual thunk
value. Since we now know that the two cells at $ptr$ are a thunk and not just an
array, we can instruct the garbage collector to replace the thunk with the
evaluated value when it encounters it. After the garbage collection, the block
at the $ptr$ location should be completely dropped, and the block in location
$x$ should just hold the value 5 directly.

This optimisation is crucial for both the memory and time efficiency of the
program. When an evaluated thunk is in memory two words of memory are wasted,
the one holding the pointer to the thunk and the \lit{Evaluated} tag. Getting
rid of these makes for more compact heaps that help the runtime efficiency of
the program. Furthermore, having the value directly available results in one
less indirection that the execution of the program has to traverse each time the
value is demanded, making accesses to it have the same performance
characteristics as if it were unboxed to begin with.

\todo[inline, size=normalsize] {
  More technical details will be inserted once all work has been completed on
  the implementation and proofs.
}
